\section{Related Work}
\label{sec:related-work}

\subsection{Educational Video Analysis}

The problem of analyzing and indexing educational video content has received considerable attention in the educational technology community. Early work focused on lecture capture systems that simply recorded and made videos available for playback \cite{copeland2009lecture}. These systems provided basic functionality but lacked intelligent content analysis.

Subsequent research introduced automatic segmentation techniques. Cui et al. \cite{cui2016automatic} developed a method for detecting slide transitions in lecture videos using color histogram analysis and template matching. Their approach achieved reasonable accuracy but struggled with subtle slide changes and animated transitions.

More recent work has leveraged deep learning for video understanding. Pan et al. \cite{pan2020educational} proposed a multimodal approach combining visual features from video frames with textual features from ASR transcripts to segment educational videos into topical units. Their method demonstrated improved segmentation quality compared to unimodal approaches.

The integration of LLMs for educational content analysis is an emerging area. Recent work by various researchers has explored using GPT-style models for generating summaries, extracting key concepts, and creating quiz questions from lecture transcripts \cite{chen2023llm}. However, these approaches typically operate on transcript text alone without leveraging visual information from slides.

\subsection{Retrieval-Augmented Generation}

RAG systems combine retrieval-based and generation-based approaches to produce grounded, factual responses. The seminal work by Lewis et al. \cite{lewis2020retrieval} introduced the RAG paradigm, demonstrating that augmenting LLMs with external knowledge retrieval significantly improves answer accuracy and reduces hallucination.

In the educational domain, several RAG-based systems have been proposed. Question-answering systems for textbooks \cite{lee2021textbook} and lecture notes \cite{wang2022lecture} have shown promising results. However, applying RAG to video content introduces unique challenges:

\begin{enumerate}
    \item \textbf{Temporal Grounding:} Answers must be linked to specific timestamps in the video, requiring precise time-range annotations in the retrieval index.
    \item \textbf{Multimodal Content:} Lecture knowledge is distributed across both spoken words and visual slides, necessitating multimodal embedding and retrieval strategies.
    \item \textbf{Sequential Context:} Video content has inherent temporal ordering that should be preserved and leveraged during retrieval and answer generation.
\end{enumerate}

Our work addresses these challenges through specialized embedding strategies that incorporate timestamp metadata and a hybrid approach that embeds both transcript text and LLM-extracted knowledge summaries.

\subsection{Video Question Answering}

Video question answering (VideoQA) is a related research area that focuses on answering questions about video content. Traditional VideoQA systems \cite{xu2017video} rely on end-to-end deep learning models trained on labeled datasets. These approaches require substantial training data and typically operate on short video clips rather than hour-long lectures.

More recent work has explored zero-shot VideoQA using pretrained LLMs with video encoders \cite{li2023zero}. While these methods show impressive generalization capabilities, they lack the fine-grained temporal grounding required for educational applications where students need to locate specific explanations within lengthy lectures.

\subsection{Knowledge Graph Construction from Text}

Automatic knowledge graph construction from textual content has been extensively studied in the NLP community. Traditional approaches use information extraction pipelines with named entity recognition, relation extraction, and entity linking \cite{hoffmann2011knowledge}. 

LLM-based approaches have recently demonstrated superior performance in extracting structured knowledge from unstructured text. Several works have explored using LLMs for concept extraction \cite{gutierrez2023llm} and ontology construction \cite{pan2023ontology}. Our fine-grained knowledge extraction task builds on this line of work, adapting it to the educational domain with structured output schemas tailored for learning objectives.
